\documentclass[pdflatex,sn-mathphys-num]{sn-jnl}

\usepackage{geometry}
\usepackage{graphicx}%
\usepackage{multirow}%
\usepackage{amsmath,amssymb,amsfonts}%
\usepackage{amsthm}%
\usepackage{mathrsfs}%
\usepackage[title]{appendix}%
\usepackage{xcolor}%
\usepackage{textcomp}%
\usepackage{manyfoot}%
\usepackage{booktabs}%
\usepackage{algorithm}%
\usepackage{algorithmicx}%
\usepackage{algpseudocode}%
\usepackage{listings}%
\usepackage{lmodern}%
\usepackage{siunitx}
\usepackage{microtype}
\usepackage{derivative}
\usepackage[version=4]{mhchem}
\usepackage{cleveref}
\usepackage{csquotes}
\usepackage{tabularx}
\usepackage{pdflscape}
\usepackage[spanish]{babel}

\renewcommand{\arraystretch}{1.2}
\sisetup{group-digits=true,
	group-separator={\,},
	separate-uncertainty
}
\AtBeginDocument{\decimalpoint}

\theoremstyle{thmstyleone}%
\newtheorem{theorem}{Teorema}%
\newtheorem{proposition}[theorem]{Proposición}%

\theoremstyle{thmstyletwo}%
\newtheorem{example}{Ejemplo}%
\newtheorem{remark}{Observación}%

\theoremstyle{thmstylethree}%
\newtheorem{definition}{Definición}%

\raggedbottom
\unnumbered% uncomment this for unnumbered level heads

\title{Revisión de Literatura sobre Nanopartículas de Plata}
\author{\fnm{Julian L.} \sur{Avila-Martinez}}
\affil{\orgdiv{Programa Académico de Física}, \orgname{Universidad Distrital
Francisco José de Caldas}, \orgaddress{\city{Bogotá D.C.}, \country{Colombia}}}
\email{jlavilam@udistrital.edu.co}
\date{\today}

\begin{document}
\maketitle

\section{Dinámica de la Literatura y Metodología de Búsqueda}
El análisis de la literatura reciente sobre nanopartículas de plata (AgNPs)
revela un sesgo profundo hacia la ingeniería aplicada y la biomedicina. La
mayoría de los estudios abordan aplicaciones macroscópicas y operativas, tales
como la administración dirigida de fármacos y la formulación de compuestos
antimicrobianos. Existe una escasez notable de investigaciones dedicadas
exclusivamente a los fenómenos físicos fundamentales de las partículas en la
escala nanométrica, como sus propiedades ópticas no lineales. Este
comportamiento del campo sugiere que las AgNPs han transitado de la teoría
física fundamental a la ingeniería aplicada.

La búsqueda que estructura la matriz de revisión obedece a los siguientes
parámetros lógicos:

\begin{itemize}
	\item Palabras clave: \enquote{Silver nanoparticles} OR \enquote{AgNPs}
	\item Rango de tiempo: 2020--2026
	\item Resultados por base de datos:
		\begin{itemize}
			\item Scopus: 6176 resultados
			\item Science Direct: 5547 resultados
			\item Springerlink: 4651 resultados
			\item Web of Science: 494 resultados
		\end{itemize}
\end{itemize}

\newpage
\newgeometry{margin=1cm}
\begin{landscape}
	\section{Matriz de Revisión de Literatura}
	\thispagestyle{empty}

	\begin{table}[htbp!]
		\caption{Artículos de revisión principales sobre AgNPs.}
		\label{tab:matriz_agnps}
		\begin{tabularx}{\linewidth}{
				>{\raggedright\arraybackslash\hsize=1.0\hsize}X
				>{\raggedright\arraybackslash\hsize=0.8\hsize}X
				>{\raggedright\arraybackslash\hsize=1.0\hsize}X
				>{\raggedright\arraybackslash\hsize=1.6\hsize}X
				>{\raggedright\arraybackslash\hsize=1.1\hsize}X
				>{\raggedright\arraybackslash\hsize=0.5\hsize}X
			}
			\toprule
			Título & Autores & Revista & Puntos principales del resumen &
			Metodología del artículo & Base \\
			\midrule
			Silver nanoparticle based polymer composites: synthesis, fabrication
			and biomedical applications - 2025 \cite{chaionSilverNanoparticleBased2025} &
			Chaion, M. H. et al. &
			Inorganic Chemistry Communications &
			Los compuestos de AgNPs integran efectos antimicrobianos y
			anticancerígenos. Se revisan métodos ecológicos frente a
			métodos químicos o físicos de síntesis, examinando su
			utilidad funcional en ingeniería de tejidos, administración de fármacos
			y biosensores. Se abordan explícitamente problemas técnicos
			como la citotoxicidad. &
			Revisión estructurada de bases de datos científicas.
			Los artículos seleccionados se categorizaron rigurosamente según sus
			métodos de síntesis, las matrices de polímeros utilizadas y el
			rendimiento biológico evaluado. &
			Science Direct \\
			\addlinespace
			A Comprehensive Review of Silver Nanoparticles (AgNPs): Synthesis
			Strategies, Toxicity Concerns, Biomedical Applications, AI-Driven
			Advancements - 2025 \cite{hosnyComprehensiveReviewSilver2025} &
			Hosny, S. et al. &
			Arabian Journal for Science and Engineering &
			Análisis de los parámetros termodinámicos que rigen la estabilidad
			de las AgNPs. Introduce la optimización algorítmica de la síntesis
			mediante inteligencia artificial para predecir citotoxicidad. &
			Evaluación sistemática de métodos de síntesis y modelos
			computacionales predictivos aplicados a la parametrización física
			en la escala nanométrica. &
			Web of Science \\
			\addlinespace
			Silver nanoparticles (AgNPs) as broad-spectrum antimicrobial agents:
			synthesis, mechanisms of action, and applications - 2026 \cite{kocSilverNanoparticlesAgNPs2026} &
			Koç, M. M. et al. &
			Journal of the Korean Physical Society &
			Evalúa la dependencia paramétrica entre los métodos de síntesis y
			la sección eficaz de interacción antimicrobiana. Describe los
			mecanismos de desintegración estructural a nivel celular. &
			Análisis correlacional de las técnicas de producción física y
			química frente a la viabilidad funcional y las restricciones de
			citotoxicidad de las nanopartículas. &
			Springerlink \\
			\addlinespace
			Silver Nanoparticles (AgNPs), Methods of Synthesis, Characterization,
			and Their Application: A Review - 2025 \cite{lasmiSilverNanoparticlesAgNPs2025} &
			Lasmi, F. et al. &
			Plasmonics &
			Examina la modulación de atributos físicos a escala atómica y la
			exigencia de caracterización estructural estricta (XRD, XPS, TEM) para
			garantizar la viabilidad funcional multimodal. &
			Revisión exhaustiva de técnicas de síntesis integrando métodos
			analíticos de caracterización espectroscópica y microscópica. &
			Web of Science \\
			\addlinespace
			Toxicological Aspects, Safety Assessment, and Green Toxicology of
			Silver Nanoparticles (AgNPs)---Critical Review: State of the Art
			- 2023 \cite{nogaToxicologicalAspectsSafety2023} &
			Noga, M. et al. &
			International Journal of Molecular Sciences &
			Análisis crítico de la dependencia funcional entre los parámetros
			fisicoquímicos de las AgNPs y su toxicidad intrínseca. Introduce la
			\enquote{toxicología verde} para mitigar daños celulares. &
			Revisión crítica de los mecanismos de toxicidad in vitro e in vivo,
			formulando directrices de seguridad basadas en la reducción de especies
			reactivas de oxígeno. &
			Scopus \\
			\addlinespace
			Nonlinear Optical Properties of Silver Nanoparticles: A Comprehensive
			Review - 2025 \cite{qayyumNonlinearOpticalProperties2025} &
			Qayyum, H. et al. &
			Laser Innovations for Research and Applications &
			Examina la polarización no lineal de las AgNPs bajo campos
			electromagnéticos intensos. Reporta fenómenos de absorción saturable
			inversa y autodesenfoque espacial derivado de un índice de refracción
			no lineal negativo. &
			Análisis teórico de la susceptibilidad óptica de tercer orden y la
			lente térmica acumulativa inducida por excitación con pulsos láser de
			femtosegundos. &
			Springerlink \\
			\addlinespace
			Silver Nanoparticles (AgNPs): Comprehensive Insights into Bio/Synthesis,
			Key Influencing Factors, Multifaceted Applications, and Toxicity---A
			2024 Update - 2025 \cite{satiSilverNanoparticlesAgNPs2025} &
			Sati, A. et al. &
			ACS Omega &
			Subraya la dependencia funcional de las propiedades ópticas y
			electrónicas de las AgNPs respecto a su método de síntesis[cite: 6].
			Argumenta la superioridad de las rutas biológicas debido a su
			estabilidad intrínseca y baja reactividad parásita[cite: 6]. &
			Análisis comparativo que contrasta la exigencia energética y la
			complejidad de purificación de enfoques físicos y químicos con la
			viabilidad de la síntesis verde mediada por fitoquímicos[cite: 6]. &
			Scopus \\
			\bottomrule
		\end{tabularx}
	\end{table}
\end{landscape}

\newpage
\newgeometry{margin=1cm}
\begin{landscape}
	\thispagestyle{empty}

	\begin{table}[htbp!]
		\caption{Continuación \cref{tab:matriz_agnps}}
		\begin{tabularx}{\linewidth}{
				>{\raggedright\arraybackslash\hsize=1.0\hsize}X
				>{\raggedright\arraybackslash\hsize=0.8\hsize}X
				>{\raggedright\arraybackslash\hsize=1.0\hsize}X
				>{\raggedright\arraybackslash\hsize=1.6\hsize}X
				>{\raggedright\arraybackslash\hsize=1.1\hsize}X
				>{\raggedright\arraybackslash\hsize=0.5\hsize}X
			}
			\toprule
			Título & Autores & Revista & Puntos principales del resumen &
			Metodología del artículo & Base \\
			\midrule
						\addlinespace
			Plant-based synthesis, characterization approaches, applications and
			toxicity of silver nanoparticles: A comprehensive review - 2024
			\cite{thomasPlantbasedSynthesisCharacterization2024} &
			Thomas, S. et al. &
			Journal of Biotechnology &
			Evalúa la síntesis verde mediada por biomoléculas vegetales que
			actúan simultáneamente como agentes reductores y estabilizadores.
			Examina la viabilidad catalítica y las restricciones de toxicidad. &
			Revisión sistemática de rutas de síntesis biogénica, integrando
			perfiles de caracterización física y evaluaciones de citotoxicidad. &
			Science Direct \\
			\bottomrule
		\end{tabularx}
	\end{table}
\end{landscape}

\restoregeometry
\newpage

\bibliography{./Silver-Nano-Particles.bib}

\end{document}
